% ===== §8 Watershed: From the Dyad to the Basin =====
\section{Watershed: From the Dyad to the Basin}
\label{sec:watershed}

\noindent
The water-virtue has been developed, so far, as an ethic of two. But water is not naturally a figure of the dyad; it is a figure of the basin, the watershed, the system of many flows that benefit and are benefited across a whole terrain. The phrase that anchors the paper says as much: water benefits the \emph{ten thousand things}, not the one other. A question therefore presses --- whether the water-virtue, worked out between two people, scales to the many: to the family, the community, the relational commons. This section opens the question and, deliberately, does not close it. The centre of gravity of this paper is the dyad, and the reasons for keeping it there are not timidity but fidelity to what has been shown; the scaling is real but its fidelity is uncertain, and the honest course is to mark the bridge and to send the crossing of it forward (\xref{sub:fissure-basin}) rather than to claim a passage the argument has not earned.

\subsection{From two drops to a watercourse}
\label{sub:watershed-scale}

The interface from which scaling would begin is already in hand. This series gave the dyad its account as the resonance of two meaning-worlds --- two drops, coupled without merging, generating a shared state irreducible to either \citep{paperxix}. The natural next object is not two drops but the watercourse: the many-bodied relational system in which value flows across more than two, and the question of its just circulation becomes the question of a commons. The resources for thinking it have been assembled in the survey. Ostrom's design principles describe how a community sustains a shared resource --- a literal watershed among her cases --- through institutions of graduated commitment and mutual monitoring that no external authority imposes \citep{ostrom1990}; the land ethic and deep ecology extend the unit of benefit from the pair to the biotic whole, the relational self finding its flourishing bound up with the flourishing of the field it belongs to \citep{leopold1949,naess1973}; and this series' own treatment of relational wealth posed the redistribution of generative value across the many, with the attendant difficulty that a system of many parties admits couplings and conflicts a dyad never faces \citep{paperxv}. ``Benefiting the ten thousand things'' is, read this way, a scaling operator: it asks the dyadic water-virtue to become the governance of a basin, the just arrangement of many flows rather than the good conduct of one coupling. The materials for that governance exist; what is not yet established is that the virtue survives the passage intact.

\subsection{The fidelity of the scaling is an open question}
\label{sub:watershed-open}

We therefore make explicit a restraint that has governed the whole paper. The dyadic water-virtue is the root, and the basin is what it might generate; but a root's extension into a larger system is not guaranteed to preserve the properties that held at its source, and there are specific reasons to expect distortion. The flesh that the formal scheme locates as the seat of intimate non-substitutability (\xref{sub:formal-flesh}) is a feature of the particular coupling; it is not obvious that it has any analogue at the scale of the many, where the parties to a commons relate under types and roles more than under the irreplaceable particularity of the beloved. The rigid lower bound that answered self-dissolution between two (\xref{sub:dialectic-bound}) may take a different form, or fail to translate, where the ``self'' at risk of being drained is a class of contributors rather than a person. And the very holonomy that carries the good-cycle criterion is, in the formal companion, sensitive to how many subjects the system contains and how the background of their relation emerges --- matters the companion explicitly leaves open, and which bear directly on whether a two-body criterion governs an n-body commons at all \citep{wan2026braid}. None of this shows that the scaling fails; it shows that the scaling is a substantive question with its own difficulties, not a corollary to be assumed. The water-virtue's passage from the dyad to the basin is accordingly handed to the first of the directed fissures (\xref{sub:fissure-basin}), as the seed of a paper in its own right --- a watershed justice that this one prepares but does not pretend to deliver. Here we hold the centre of gravity where the argument has actually done its work: between two.
